\section{}

\paragraph{1013年}\eventanchor{XIA-1013-REGNAL-YEAR}{XIA-1013-REGNAL-YEAR}　党项李氏以李德明在位第9年作为诏令、赋役与外交文书的纪年框架。账本列出《宋史》卷485〈夏国传〉；《辽史》《金史》相关本纪；黑水城西夏文文书与考古报告作为来源线索；该材料可核对时间与制度称谓，但有关执行范围、群体经验、军数或后果，仍须与地方文书、物质遗存及相关政权记录互证。

\paragraph{1014年}\eventanchor{XIA-1014-REGNAL-YEAR}{XIA-1014-REGNAL-YEAR}　党项李氏以李德明在位第10年作为诏令、赋役与外交文书的纪年框架。账本列出《宋史》卷485〈夏国传〉；《辽史》《金史》相关本纪；黑水城西夏文文书与考古报告作为来源线索；该材料可核对时间与制度称谓，但有关执行范围、群体经验、军数或后果，仍须与地方文书、物质遗存及相关政权记录互证。

\paragraph{1015年}\eventanchor{XIA-1015-REGNAL-YEAR}{XIA-1015-REGNAL-YEAR}　党项李氏以李德明在位第11年作为诏令、赋役与外交文书的纪年框架。账本列出《宋史》卷485〈夏国传〉；《辽史》《金史》相关本纪；黑水城西夏文文书与考古报告作为来源线索；该材料可核对时间与制度称谓，但有关执行范围、群体经验、军数或后果，仍须与地方文书、物质遗存及相关政权记录互证。

\paragraph{1016年}\eventanchor{XIA-1016-REGNAL-YEAR}{XIA-1016-REGNAL-YEAR}　党项李氏以李德明在位第12年作为诏令、赋役与外交文书的纪年框架。账本列出《宋史》卷485〈夏国传〉；《辽史》《金史》相关本纪；黑水城西夏文文书与考古报告作为来源线索；该材料可核对时间与制度称谓，但有关执行范围、群体经验、军数或后果，仍须与地方文书、物质遗存及相关政权记录互证。

\paragraph{1017年}\eventanchor{XIA-1017-REGNAL-YEAR}{XIA-1017-REGNAL-YEAR}　党项李氏以李德明在位第13年作为诏令、赋役与外交文书的纪年框架。账本列出《宋史》卷485〈夏国传〉；《辽史》《金史》相关本纪；黑水城西夏文文书与考古报告作为来源线索；该材料可核对时间与制度称谓，但有关执行范围、群体经验、军数或后果，仍须与地方文书、物质遗存及相关政权记录互证。

\paragraph{1018年}\eventanchor{XIA-1018-REGNAL-YEAR}{XIA-1018-REGNAL-YEAR}　党项李氏以李德明在位第14年作为诏令、赋役与外交文书的纪年框架。账本列出《宋史》卷485〈夏国传〉；《辽史》《金史》相关本纪；黑水城西夏文文书与考古报告作为来源线索；该材料可核对时间与制度称谓，但有关执行范围、群体经验、军数或后果，仍须与地方文书、物质遗存及相关政权记录互证。

\paragraph{1019年}\eventanchor{XIA-1019-REGNAL-YEAR}{XIA-1019-REGNAL-YEAR}　党项李氏以李德明在位第15年作为诏令、赋役与外交文书的纪年框架。账本列出《宋史》卷485〈夏国传〉；《辽史》《金史》相关本纪；黑水城西夏文文书与考古报告作为来源线索；该材料可核对时间与制度称谓，但有关执行范围、群体经验、军数或后果，仍须与地方文书、物质遗存及相关政权记录互证。

\paragraph{1020年}\eventanchor{XIA-1020-REGNAL-YEAR}{XIA-1020-REGNAL-YEAR}　党项李氏以李德明在位第16年作为诏令、赋役与外交文书的纪年框架。账本列出《宋史》卷485〈夏国传〉；《辽史》《金史》相关本纪；黑水城西夏文文书与考古报告作为来源线索；该材料可核对时间与制度称谓，但有关执行范围、群体经验、军数或后果，仍须与地方文书、物质遗存及相关政权记录互证。

\paragraph{1021年}\eventanchor{XIA-1021-REGNAL-YEAR}{XIA-1021-REGNAL-YEAR}　党项李氏以李德明在位第17年作为诏令、赋役与外交文书的纪年框架。账本列出《宋史》卷485〈夏国传〉；《辽史》《金史》相关本纪；黑水城西夏文文书与考古报告作为来源线索；该材料可核对时间与制度称谓，但有关执行范围、群体经验、军数或后果，仍须与地方文书、物质遗存及相关政权记录互证。

\paragraph{1022年}\eventanchor{XIA-1022-REGNAL-YEAR}{XIA-1022-REGNAL-YEAR}　党项李氏以李德明在位第18年作为诏令、赋役与外交文书的纪年框架。账本列出《宋史》卷485〈夏国传〉；《辽史》《金史》相关本纪；黑水城西夏文文书与考古报告作为来源线索；该材料可核对时间与制度称谓，但有关执行范围、群体经验、军数或后果，仍须与地方文书、物质遗存及相关政权记录互证。

\paragraph{1023年}\eventanchor{XIA-1023-REGNAL-YEAR}{XIA-1023-REGNAL-YEAR}　党项李氏以李德明在位第19年作为诏令、赋役与外交文书的纪年框架。账本列出《宋史》卷485〈夏国传〉；《辽史》《金史》相关本纪；黑水城西夏文文书与考古报告作为来源线索；该材料可核对时间与制度称谓，但有关执行范围、群体经验、军数或后果，仍须与地方文书、物质遗存及相关政权记录互证。

\paragraph{1024年}\eventanchor{XIA-1024-REGNAL-YEAR}{XIA-1024-REGNAL-YEAR}　党项李氏以李德明在位第20年作为诏令、赋役与外交文书的纪年框架。账本列出《宋史》卷485〈夏国传〉；《辽史》《金史》相关本纪；黑水城西夏文文书与考古报告作为来源线索；该材料可核对时间与制度称谓，但有关执行范围、群体经验、军数或后果，仍须与地方文书、物质遗存及相关政权记录互证。

\paragraph{1025年}\eventanchor{XIA-1025-REGNAL-YEAR}{XIA-1025-REGNAL-YEAR}　党项李氏以李德明在位第21年作为诏令、赋役与外交文书的纪年框架。账本列出《宋史》卷485〈夏国传〉；《辽史》《金史》相关本纪；黑水城西夏文文书与考古报告作为来源线索；该材料可核对时间与制度称谓，但有关执行范围、群体经验、军数或后果，仍须与地方文书、物质遗存及相关政权记录互证。

\paragraph{1026年}\eventanchor{XIA-1026-REGNAL-YEAR}{XIA-1026-REGNAL-YEAR}　党项李氏以李德明在位第22年作为诏令、赋役与外交文书的纪年框架。账本列出《宋史》卷485〈夏国传〉；《辽史》《金史》相关本纪；黑水城西夏文文书与考古报告作为来源线索；该材料可核对时间与制度称谓，但有关执行范围、群体经验、军数或后果，仍须与地方文书、物质遗存及相关政权记录互证。

\paragraph{1027年}\eventanchor{XIA-1027-REGNAL-YEAR}{XIA-1027-REGNAL-YEAR}　党项李氏以李德明在位第23年作为诏令、赋役与外交文书的纪年框架。账本列出《宋史》卷485〈夏国传〉；《辽史》《金史》相关本纪；黑水城西夏文文书与考古报告作为来源线索；该材料可核对时间与制度称谓，但有关执行范围、群体经验、军数或后果，仍须与地方文书、物质遗存及相关政权记录互证。

\paragraph{1028年}\eventanchor{XIA-1028-REGNAL-YEAR}{XIA-1028-REGNAL-YEAR}　党项李氏以李德明在位第24年作为诏令、赋役与外交文书的纪年框架。账本列出《宋史》卷485〈夏国传〉；《辽史》《金史》相关本纪；黑水城西夏文文书与考古报告作为来源线索；该材料可核对时间与制度称谓，但有关执行范围、群体经验、军数或后果，仍须与地方文书、物质遗存及相关政权记录互证。

\paragraph{1029年}\eventanchor{XIA-1029-REGNAL-YEAR}{XIA-1029-REGNAL-YEAR}　党项李氏以李德明在位第25年作为诏令、赋役与外交文书的纪年框架。账本列出《宋史》卷485〈夏国传〉；《辽史》《金史》相关本纪；黑水城西夏文文书与考古报告作为来源线索；该材料可核对时间与制度称谓，但有关执行范围、群体经验、军数或后果，仍须与地方文书、物质遗存及相关政权记录互证。

\paragraph{1030年}\eventanchor{XIA-1030-REGNAL-YEAR}{XIA-1030-REGNAL-YEAR}　党项李氏以李德明在位第26年作为诏令、赋役与外交文书的纪年框架。账本列出《宋史》卷485〈夏国传〉；《辽史》《金史》相关本纪；黑水城西夏文文书与考古报告作为来源线索；该材料可核对时间与制度称谓，但有关执行范围、群体经验、军数或后果，仍须与地方文书、物质遗存及相关政权记录互证。

\paragraph{1031年}\eventanchor{XIA-1031-REGNAL-YEAR}{XIA-1031-REGNAL-YEAR}　党项李氏以李德明在位第27年作为诏令、赋役与外交文书的纪年框架。账本列出《宋史》卷485〈夏国传〉；《辽史》《金史》相关本纪；黑水城西夏文文书与考古报告作为来源线索；该材料可核对时间与制度称谓，但有关执行范围、群体经验、军数或后果，仍须与地方文书、物质遗存及相关政权记录互证。

\paragraph{1032年}\eventanchor{XIA-1032-EVENT-053}{XIA-1032-EVENT-053}　李德明去世，李元昊继位。账本列出《宋史》卷485〈夏国传〉；《辽史》《金史》相关本纪；黑水城西夏文文书与考古报告作为来源线索；该材料可核对时间与制度称谓，但有关执行范围、群体经验、军数或后果，仍须与地方文书、物质遗存及相关政权记录互证。

\paragraph{1032年}\eventanchor{XIA-1032-REGNAL-YEAR}{XIA-1032-REGNAL-YEAR}　党项李氏以李德明在位第28年作为诏令、赋役与外交文书的纪年框架。账本列出《宋史》卷485〈夏国传〉；《辽史》《金史》相关本纪；黑水城西夏文文书与考古报告作为来源线索；该材料可核对时间与制度称谓，但有关执行范围、群体经验、军数或后果，仍须与地方文书、物质遗存及相关政权记录互证。

\paragraph{1033年}\eventanchor{XIA-1033-REGNAL-YEAR}{XIA-1033-REGNAL-YEAR}　西夏以李元昊在位第1年作为诏令、赋役与外交文书的纪年框架。账本列出《宋史》卷485〈夏国传〉；《辽史》《金史》相关本纪；黑水城西夏文文书与考古报告作为来源线索；该材料可核对时间与制度称谓，但有关执行范围、群体经验、军数或后果，仍须与地方文书、物质遗存及相关政权记录互证。

\paragraph{1034年}\eventanchor{XIA-1034-EVENT-054}{XIA-1034-EVENT-054}　李元昊改国号、年号并改革王权象征。账本列出《宋史》卷485〈夏国传〉；《辽史》《金史》相关本纪；黑水城西夏文文书与考古报告作为来源线索；该材料可核对时间与制度称谓，但有关执行范围、群体经验、军数或后果，仍须与地方文书、物质遗存及相关政权记录互证。

\paragraph{1034年}\eventanchor{XIA-1034-REGNAL-YEAR}{XIA-1034-REGNAL-YEAR}　西夏以李元昊在位第2年作为诏令、赋役与外交文书的纪年框架。账本列出《宋史》卷485〈夏国传〉；《辽史》《金史》相关本纪；黑水城西夏文文书与考古报告作为来源线索；该材料可核对时间与制度称谓，但有关执行范围、群体经验、军数或后果，仍须与地方文书、物质遗存及相关政权记录互证。

\paragraph{1035年}\eventanchor{XIA-1035-REGNAL-YEAR}{XIA-1035-REGNAL-YEAR}　西夏以李元昊在位第3年作为诏令、赋役与外交文书的纪年框架。账本列出《宋史》卷485〈夏国传〉；《辽史》《金史》相关本纪；黑水城西夏文文书与考古报告作为来源线索；该材料可核对时间与制度称谓，但有关执行范围、群体经验、军数或后果，仍须与地方文书、物质遗存及相关政权记录互证。

\paragraph{1036年}\eventanchor{XIA-1036-EVENT-055}{XIA-1036-EVENT-055}　党项军向甘州、凉州方向扩张，控制河西要点。账本列出《宋史》卷485〈夏国传〉；《辽史》《金史》相关本纪；黑水城西夏文文书与考古报告作为来源线索；该材料可核对时间与制度称谓，但有关执行范围、群体经验、军数或后果，仍须与地方文书、物质遗存及相关政权记录互证。

\paragraph{1036年}\eventanchor{XIA-1036-REGNAL-YEAR}{XIA-1036-REGNAL-YEAR}　西夏以李元昊在位第4年作为诏令、赋役与外交文书的纪年框架。账本列出《宋史》卷485〈夏国传〉；《辽史》《金史》相关本纪；黑水城西夏文文书与考古报告作为来源线索；该材料可核对时间与制度称谓，但有关执行范围、群体经验、军数或后果，仍须与地方文书、物质遗存及相关政权记录互证。

\paragraph{1037年}\eventanchor{XIA-1037-REGNAL-YEAR}{XIA-1037-REGNAL-YEAR}　西夏以李元昊在位第5年作为诏令、赋役与外交文书的纪年框架。账本列出《宋史》卷485〈夏国传〉；《辽史》《金史》相关本纪；黑水城西夏文文书与考古报告作为来源线索；该材料可核对时间与制度称谓，但有关执行范围、群体经验、军数或后果，仍须与地方文书、物质遗存及相关政权记录互证。

\paragraph{1038年}\eventanchor{XIA-1038-EVENT-056}{XIA-1038-EVENT-056}　李元昊称帝，西夏正式建国。账本列出《宋史》卷485〈夏国传〉；《辽史》《金史》相关本纪；黑水城西夏文文书与考古报告作为来源线索；该材料可核对时间与制度称谓，但有关执行范围、群体经验、军数或后果，仍须与地方文书、物质遗存及相关政权记录互证。

\paragraph{1038年}\eventanchor{XIA-1038-REGNAL-YEAR}{XIA-1038-REGNAL-YEAR}　西夏以李元昊在位第6年作为诏令、赋役与外交文书的纪年框架。账本列出《宋史》卷485〈夏国传〉；《辽史》《金史》相关本纪；黑水城西夏文文书与考古报告作为来源线索；该材料可核对时间与制度称谓，但有关执行范围、群体经验、军数或后果，仍须与地方文书、物质遗存及相关政权记录互证。

\paragraph{1039年}\eventanchor{XIA-1039-REGNAL-YEAR}{XIA-1039-REGNAL-YEAR}　西夏以李元昊在位第7年作为诏令、赋役与外交文书的纪年框架。账本列出《宋史》卷485〈夏国传〉；《辽史》《金史》相关本纪；黑水城西夏文文书与考古报告作为来源线索；该材料可核对时间与制度称谓，但有关执行范围、群体经验、军数或后果，仍须与地方文书、物质遗存及相关政权记录互证。

\paragraph{1040年}\eventanchor{XIA-1040-EVENT-057}{XIA-1040-EVENT-057}　三川口之战中宋军失利。账本列出《宋史》卷485〈夏国传〉；《辽史》《金史》相关本纪；黑水城西夏文文书与考古报告作为来源线索；该材料可核对时间与制度称谓，但有关执行范围、群体经验、军数或后果，仍须与地方文书、物质遗存及相关政权记录互证。

\paragraph{1040年}\eventanchor{XIA-1040-REGNAL-YEAR}{XIA-1040-REGNAL-YEAR}　西夏以李元昊在位第8年作为诏令、赋役与外交文书的纪年框架。账本列出《宋史》卷485〈夏国传〉；《辽史》《金史》相关本纪；黑水城西夏文文书与考古报告作为来源线索；该材料可核对时间与制度称谓，但有关执行范围、群体经验、军数或后果，仍须与地方文书、物质遗存及相关政权记录互证。

\paragraph{1041年}\eventanchor{XIA-1041-EVENT-058}{XIA-1041-EVENT-058}　好水川之战使宋军再遭挫败。账本列出《宋史》卷485〈夏国传〉；《辽史》《金史》相关本纪；黑水城西夏文文书与考古报告作为来源线索；该材料可核对时间与制度称谓，但有关执行范围、群体经验、军数或后果，仍须与地方文书、物质遗存及相关政权记录互证。

\paragraph{1041年}\eventanchor{XIA-1041-REGNAL-YEAR}{XIA-1041-REGNAL-YEAR}　西夏以李元昊在位第9年作为诏令、赋役与外交文书的纪年框架。账本列出《宋史》卷485〈夏国传〉；《辽史》《金史》相关本纪；黑水城西夏文文书与考古报告作为来源线索；该材料可核对时间与制度称谓，但有关执行范围、群体经验、军数或后果，仍须与地方文书、物质遗存及相关政权记录互证。

\paragraph{1042年}\eventanchor{XIA-1042-EVENT-059}{XIA-1042-EVENT-059}　定川寨之战加剧宋西北防务危机。账本列出《宋史》卷485〈夏国传〉；《辽史》《金史》相关本纪；黑水城西夏文文书与考古报告作为来源线索；该材料可核对时间与制度称谓，但有关执行范围、群体经验、军数或后果，仍须与地方文书、物质遗存及相关政权记录互证。

\paragraph{1042年}\eventanchor{XIA-1042-REGNAL-YEAR}{XIA-1042-REGNAL-YEAR}　西夏以李元昊在位第10年作为诏令、赋役与外交文书的纪年框架。账本列出《宋史》卷485〈夏国传〉；《辽史》《金史》相关本纪；黑水城西夏文文书与考古报告作为来源线索；该材料可核对时间与制度称谓，但有关执行范围、群体经验、军数或后果，仍须与地方文书、物质遗存及相关政权记录互证。

\paragraph{1043年}\eventanchor{XIA-1043-REGNAL-YEAR}{XIA-1043-REGNAL-YEAR}　西夏以李元昊在位第11年作为诏令、赋役与外交文书的纪年框架。账本列出《宋史》卷485〈夏国传〉；《辽史》《金史》相关本纪；黑水城西夏文文书与考古报告作为来源线索；该材料可核对时间与制度称谓，但有关执行范围、群体经验、军数或后果，仍须与地方文书、物质遗存及相关政权记录互证。

\paragraph{1044年}\eventanchor{XIA-1044-EVENT-060}{XIA-1044-EVENT-060}　宋夏达成庆历和议。账本列出《宋史》卷485〈夏国传〉；《辽史》《金史》相关本纪；黑水城西夏文文书与考古报告作为来源线索；该材料可核对时间与制度称谓，但有关执行范围、群体经验、军数或后果，仍须与地方文书、物质遗存及相关政权记录互证。

\paragraph{1044年}\eventanchor{XIA-1044-REGNAL-YEAR}{XIA-1044-REGNAL-YEAR}　西夏以李元昊在位第12年作为诏令、赋役与外交文书的纪年框架。账本列出《宋史》卷485〈夏国传〉；《辽史》《金史》相关本纪；黑水城西夏文文书与考古报告作为来源线索；该材料可核对时间与制度称谓，但有关执行范围、群体经验、军数或后果，仍须与地方文书、物质遗存及相关政权记录互证。

\paragraph{1045年}\eventanchor{XIA-1045-REGNAL-YEAR}{XIA-1045-REGNAL-YEAR}　西夏以李元昊在位第13年作为诏令、赋役与外交文书的纪年框架。账本列出《宋史》卷485〈夏国传〉；《辽史》《金史》相关本纪；黑水城西夏文文书与考古报告作为来源线索；该材料可核对时间与制度称谓，但有关执行范围、群体经验、军数或后果，仍须与地方文书、物质遗存及相关政权记录互证。

\paragraph{1046年}\eventanchor{XIA-1046-REGNAL-YEAR}{XIA-1046-REGNAL-YEAR}　西夏以李元昊在位第14年作为诏令、赋役与外交文书的纪年框架。账本列出《宋史》卷485〈夏国传〉；《辽史》《金史》相关本纪；黑水城西夏文文书与考古报告作为来源线索；该材料可核对时间与制度称谓，但有关执行范围、群体经验、军数或后果，仍须与地方文书、物质遗存及相关政权记录互证。

\paragraph{1047年}\eventanchor{XIA-1047-REGNAL-YEAR}{XIA-1047-REGNAL-YEAR}　西夏以李元昊在位第15年作为诏令、赋役与外交文书的纪年框架。账本列出《宋史》卷485〈夏国传〉；《辽史》《金史》相关本纪；黑水城西夏文文书与考古报告作为来源线索；该材料可核对时间与制度称谓，但有关执行范围、群体经验、军数或后果，仍须与地方文书、物质遗存及相关政权记录互证。

\paragraph{1048年}\eventanchor{XIA-1048-EVENT-061}{XIA-1048-EVENT-061}　李元昊遇害，幼主继位并出现后族、权臣主政。账本列出《宋史》卷485〈夏国传〉；《辽史》《金史》相关本纪；黑水城西夏文文书与考古报告作为来源线索；该材料可核对时间与制度称谓，但有关执行范围、群体经验、军数或后果，仍须与地方文书、物质遗存及相关政权记录互证。

\paragraph{1048年}\eventanchor{XIA-1048-REGNAL-YEAR}{XIA-1048-REGNAL-YEAR}　西夏以李元昊在位第16年作为诏令、赋役与外交文书的纪年框架。账本列出《宋史》卷485〈夏国传〉；《辽史》《金史》相关本纪；黑水城西夏文文书与考古报告作为来源线索；该材料可核对时间与制度称谓，但有关执行范围、群体经验、军数或后果，仍须与地方文书、物质遗存及相关政权记录互证。

\paragraph{1049年}\eventanchor{XIA-1049-EVENT-062}{XIA-1049-EVENT-062}　李谅祚即位。账本列出《宋史》卷485〈夏国传〉；《辽史》《金史》相关本纪；黑水城西夏文文书与考古报告作为来源线索；该材料可核对时间与制度称谓，但有关执行范围、群体经验、军数或后果，仍须与地方文书、物质遗存及相关政权记录互证。

\paragraph{1049年}\eventanchor{XIA-1049-REGNAL-YEAR}{XIA-1049-REGNAL-YEAR}　西夏以李谅祚在位第1年作为诏令、赋役与外交文书的纪年框架。账本列出《宋史》卷485〈夏国传〉；《辽史》《金史》相关本纪；黑水城西夏文文书与考古报告作为来源线索；该材料可核对时间与制度称谓，但有关执行范围、群体经验、军数或后果，仍须与地方文书、物质遗存及相关政权记录互证。

\paragraph{1050年}\eventanchor{XIA-1050-REGNAL-YEAR}{XIA-1050-REGNAL-YEAR}　西夏以李谅祚在位第2年作为诏令、赋役与外交文书的纪年框架。账本列出《宋史》卷485〈夏国传〉；《辽史》《金史》相关本纪；黑水城西夏文文书与考古报告作为来源线索；该材料可核对时间与制度称谓，但有关执行范围、群体经验、军数或后果，仍须与地方文书、物质遗存及相关政权记录互证。

\paragraph{1051年}\eventanchor{XIA-1051-REGNAL-YEAR}{XIA-1051-REGNAL-YEAR}　西夏以李谅祚在位第3年作为诏令、赋役与外交文书的纪年框架。账本列出《宋史》卷485〈夏国传〉；《辽史》《金史》相关本纪；黑水城西夏文文书与考古报告作为来源线索；该材料可核对时间与制度称谓，但有关执行范围、群体经验、军数或后果，仍须与地方文书、物质遗存及相关政权记录互证。

\paragraph{1052年}\eventanchor{XIA-1052-REGNAL-YEAR}{XIA-1052-REGNAL-YEAR}　西夏以李谅祚在位第4年作为诏令、赋役与外交文书的纪年框架。账本列出《宋史》卷485〈夏国传〉；《辽史》《金史》相关本纪；黑水城西夏文文书与考古报告作为来源线索；该材料可核对时间与制度称谓，但有关执行范围、群体经验、军数或后果，仍须与地方文书、物质遗存及相关政权记录互证。

